\clearpage\SourceLeaf{116}
\pdfbookmark[1]{Göttingische gelehrte Anzeigen. 2. März 1816.}{anzeige1816}
\begin{center}
\rule{0.45\linewidth}{0.4pt}\par
\textbf{Göttingische gelehrte Anzeigen.\quad 2. März 1816.}\par
\rule{0.45\linewidth}{0.4pt}
\end{center}

Nachdem die Vorlesung von Herrn Professor \textsc{Gauss}, über die in diesen Anzeigen im vorigen Jahr berichtet wurde, bereits gedruckt war, hatte er bei der weiteren Beschäftigung mit demselben Gegenstand das Glück, dasselbe Ziel auf einem ganz neuen Weg zu erreichen. Er hat darüber am 30. Januar der Königlichen Gesellschaft eine kurze Abhandlung mit folgendem Titel vorgelegt:

\begin{center}
\emph{Dritter Beweis des Satzes über die Zerlegbarkeit von Polynomen in reelle Faktoren.}
\end{center}

Die beiden früheren Beweise des wichtigsten Satzes in der Theorie der Gleichungen unterschieden sich dadurch, dass der erste sehr kurz und einfach war, aber zum Teil auf geometrischen Betrachtungen beruhte, während der andere rein analytisch, aber viel komplizierter war, sodass sich in dem engen Raum dieser Blätter kein ausreichender Auszug daraus geben ließ. Der jetzige dritte Beweis beruht auf völlig anderen Prinzipien und ist ebenfalls rein analytisch. Er übertrifft aber selbst den ersten so sehr an Einfachheit und Kürze, dass wir hier ohne Weiteres alles Wesentliche auf wenigen Seiten mitteilen können.

Es bezeichne $X$ die algebraische Funktion der Unbestimmten $x$,
\[
x^m+Ax^{m-1}+Bx^{m-2}+Cx^{m-3}+\mathrm{usw.}
\]
\SourceLeaf{117}
wobei die Koeffizienten $A,B,C$ usw. bestimmte reelle Größen sind. Ferner seien $r$ und $\varphi$ zwei neue Unbestimmte, und es sei
\[
\begin{aligned}
&r^m\cos m\varphi+Ar^{m-1}\cos(m-1)\varphi
\\
&\quad+Br^{m-2}\cos(m-2)\varphi+Cr^{m-3}\cos(m-3)\varphi
+\mathrm{usw.}=t\\
&r^m\sin m\varphi+Ar^{m-1}\sin(m-1)\varphi
\\
&\quad+Br^{m-2}\sin(m-2)\varphi+Cr^{m-3}\sin(m-3)\varphi
+\mathrm{usw.}=u\\
&mr^m\cos m\varphi+(m-1)Ar^{m-1}\cos(m-1)\varphi
\\
&\quad+(m-2)Br^{m-2}\cos(m-2)\varphi+(m-3)Cr^{m-3}\cos(m-3)\varphi
+\mathrm{usw.}=t'\\
&mr^m\sin m\varphi+(m-1)Ar^{m-1}\sin(m-1)\varphi
\\
&\quad+(m-2)Br^{m-2}\sin(m-2)\varphi+(m-3)Cr^{m-3}\sin(m-3)\varphi
+\mathrm{usw.}=u'\\
&mmr^m\cos m\varphi+(m-1)^2Ar^{m-1}\cos(m-1)\varphi
\\
&\quad+(m-2)^2Br^{m-2}\cos(m-2)\varphi\\
&\quad+(m-3)^2Cr^{m-3}\cos(m-3)\varphi
+\mathrm{usw.}=t''\\
&mmr^m\sin m\varphi+(m-1)^2Ar^{m-1}\sin(m-1)\varphi
\\
&\quad+(m-2)^2Br^{m-2}\sin(m-2)\varphi\\
&\quad+(m-3)^2Cr^{m-3}\sin(m-3)\varphi
+\mathrm{usw.}=u''\\
&\frac{(tt+uu)(tt''+uu'')+(tu'-ut')^2-(tt'+uu')^2}
{r(tt+uu)^2}=y
\end{aligned}
\]

Der Faktor $r$ kann offenbar aus dem Nenner des Ausdrucks für $y$ entfernt werden, da $t',u',t'',u''$ durch $r$ teilbar sind.

Ferner sei $R$ eine beliebige bestimmte positive Größe, mit der einzigen Einschränkung, dass sie größer ist als die größte der folgenden Größen:
\[
mA\surd2,\quad\surd(mB\surd2),\quad
{}^3\!\surd(mC\surd2),\quad{}^4\!\surd(mD\surd2)\ \mathrm{usw.}
\]
wobei von den Vorzeichen von $A,B,C$ usw. abgesehen wird oder alle als positiv betrachtet werden.

Schließlich seien noch folgende Bezeichnungen eingeführt:
\[
\begin{aligned}
&R^m\cos45^0+AR^{m-1}\cos(45^0+\varphi)
\\
&\quad+BR^{m-2}\cos(45^0+2\varphi)+CR^{m-3}\cos(45^0+3\varphi)
+\mathrm{usw.}=T\\
&R^m\sin45^0+AR^{m-1}\sin(45^0+\varphi)
\\
&\quad+BR^{m-2}\sin(45^0+2\varphi)+CR^{m-3}\sin(45^0+3\varphi)
+\mathrm{usw.}=U\\
&mR^m\cos45^0+(m-1)AR^{m-1}\cos(45^0+\varphi)
\\
&\quad+(m-2)BR^{m-2}\cos(45^0+2\varphi)
\\
&\quad+(m-3)CR^{m-3}\cos(45^0+3\varphi)+\mathrm{usw.}=T'\\
&mR^m\sin45^0+(m-1)AR^{m-1}\sin(45^0+\varphi)
\\
&\quad+(m-2)BR^{m-2}\sin(45^0+2\varphi)
\\
&\quad+(m-3)CR^{m-3}\sin(45^0+3\varphi)+\mathrm{usw.}=U'
\end{aligned}
\]

Nach all diesen Vorbereitungen, die wir zur leichteren Übersicht zusammengestellt haben, ergeben sich leicht die folgenden Schlüsse:

I. Die Größe $T$ ist notwendig positiv, welchen Wert auch immer $\varphi$ erhält. Dies ist leicht zu erkennen, wenn $T$ in die folgende Form gebracht wird:
\[
\begin{aligned}
&\frac{R^{m-1}}{m\surd2}(R+mA\surd2.\cos(45^{\mathrm{a}}+\varphi))\\
&\quad+\frac{R^{m-2}}{m\surd2}(RR+mB\surd2.\cos(45^0+2\varphi))\\
&\quad+\frac{R^{m-3}}{m\surd2}(R^3+mC\surd2.\cos(45^0+3\varphi))\\
&\quad+\frac{R^{m-4}}{m\surd2}(R^4+mD\surd2.\cos(45^0+4\varphi))\\
&\quad+\mathrm{usw.}
\end{aligned}
\]
\SourceLeaf{118}
denn jeder Teil wird für sich genommen notwendig positiv. Auf ähnliche Weise lässt sich leicht beweisen, dass für jeden Wert von $\varphi$ auch $U,T',U'$ notwendig positiv werden.

II. Für $r=R$ nehmen die Funktionen $t,u,t',u'$ der Reihe nach folgende Werte an:
\[
\begin{gathered}
T\cos(45^0+m\varphi)+U\sin(45^0+m\varphi)\\
T\sin(45^0+m\varphi)-U\cos(45^0+m\varphi)\\
T'\cos(45^0+m\varphi)+U'\sin(45^0+m\varphi)\\
T'\sin(45^0+m\varphi)-U'\cos(45^0+m\varphi)
\end{gathered}
\]
Daraus folgt, dass für $r=R$, $tt+uu=TT+UU$, $tt'+uu'=TT'+UU'$, also beide notwendig positiv werden.

III. Somit können für keinen Wert von $r$, der größer ist als jede der Größen $mA\surd2$, $\surd(mB\surd2)$, ${}^3\!\surd(mC\surd2)$ usw., zugleich $t=0$ und $u=0$ gelten.

Der Hauptsatz lautet nun:

\emph{Innerhalb der Grenzen $r=0$, $r=R$ und $\varphi=0$, $\varphi=360^0$ (einschließlich) gibt es gewiss Werte für $r$ und $\varphi$, die zugleich $t=0$ und $u=0$ ergeben.}

Der Beweis wird folgendermaßen geführt. Wird angenommen, der Satz sei nicht wahr, so folgt, dass $tt+uu$ für alle Werte von $r$ und $\varphi$ zwischen den angegebenen Grenzen stets positiv und folglich $y$ stets \emph{endlich} sein muss. Betrachten wir nun den Wert des Doppelintegrals
\[
\iint y\,\mathrm{d}r\,\mathrm{d}\varphi
\]
von $r=0$ bis $r=R$ und von $\varphi=0$ bis $\varphi=360^0$, der also eine ganz bestimmte endliche Größe haben wird. Dieser mit $\Omega$ zu bezeichnende Wert lässt sich auf zwei Arten finden: durch Integration zuerst nach $\varphi$ und dann nach $r$ oder in umgekehrter Reihenfolge. Beide Wege müssen notwendig zum selben Ergebnis führen.

Nun gilt aber, wenn $r$ als konstant betrachtet wird, unbestimmt
\[
\int y\,\mathrm{d}\varphi=\frac{tu'-ut'}{r(tt+uu)}
\]
wie sich durch Differentiation leicht nachprüfen lässt. Eine konstante Größe
\SourceLeaf{119}
ist bei der Integration nicht hinzuzufügen, sofern diese bei $\varphi=0$ beginnen soll, denn für diesen Wert gilt offenbar
\[
\frac{tu'-ut'}{tt+uu}=0
\]
Da dasselbe für $\varphi=360^0$ gilt, ist der Wert von $\int y\,d\varphi$, von $\varphi=0$ bis $\varphi=360^0$ genommen, $=0$, und zwar bei jedem Wert von $r$. Daraus schließen wir also $\Omega=0$.

Andererseits gilt, wenn $\varphi$ als konstant betrachtet wird, unbestimmt
\[
\int y\,dr=\frac{tt'+uu'}{tt+uu}
\]
wie ebenfalls durch Differentiation nach $r$ leicht bestätigt wird. Auch hier ist keine konstante Größe hinzuzufügen, sofern die Integration bei $r=0$ beginnen soll. Das Integral $\int y\,dr$, von $r=0$ bis $r=R$ genommen, wird folglich nach II
\[
=\frac{TT'+UU'}{TT+UU}
\]
also gewiss positiv für jeden Wert von $\varphi$. Daher wird auch $\Omega$, also der Wert des Integrals
\[
\int\frac{TT'+UU'}{TT+UU}\,d\varphi
\]
von $\varphi=0$ bis $\varphi=360^0$ genommen, notwendig positiv, was dem ersten Ergebnis widerspricht. Die Annahme war folglich unzulässig, und damit ist die Wahrheit des Satzes selbst bewiesen.

Dass für alle Werte von $r$ und $\varphi$, die zugleich $t=0$ und $u=0$ ergeben,
\[
r(\cos\varphi+\sin\varphi\surd(-1))\quad\hbox{und}\quad
r(\cos\varphi-\sin\varphi\surd(-1))
\]
Wurzeln der Gleichung $X=0$ sind oder dass die Funktion $X$ dann entweder den Faktor zweiten Grades
\[
xx-2r\cos\varphi.x+rr
\]
(wenn weder $r$ noch $\sin\varphi$ verschwindet) oder den Faktor ersten Grades
\[
x\mp r
\]
\SourceLeaf{120}
(wenn entweder $\sin\varphi=0$ und somit $\cos\varphi=\pm1$ oder $r=0$ gilt) enthält, ist hinreichend bekannt und braucht hier nicht weiter ausgeführt zu werden. Was gerade von $X$ bewiesen wurde, gilt anschließend wieder für den Quotienten, wenn $X$ durch jenen Faktor dividiert ist, und so fort. Daraus ergibt sich die vollständige Zerlegbarkeit der Funktion $X$ in solche Faktoren.

Mit dieser kurzen Folge von Schlüssen ist das eigentliche Ziel vollständig erreicht. Nicht gefordert, aber doch gewünscht werden könnte noch eine Erläuterung, inwiefern der obige Widerspruch, der unvermeidlich war, wenn der Satz als unwahr angesehen wurde, entfällt, wenn von der Wahrheit des Satzes ausgegangen wird. Der Verfasser hat dazu einige Hinweise gegeben, die hier nur kurz berührt werden können. Ist der Satz wahr, so sind nicht mehr alle Werte von $y$ endlich. Folglich ist es nicht mehr ohne Weiteres erlaubt, $\iint y\,dr\,d\varphi$ als eine wirkliche bestimmte Größe anzusehen. Es darf daher nicht verwundern, dass der bloße blinde Mechanismus des Kalküls auf verschiedenen Wegen widersprechende Ergebnisse liefert. Die Analysis hilft sich sehr häufig so, dass sie auf Fragen, die ihr ohne Einschränkung vorgelegt werden, obwohl sie in gewissen Fällen unsinnig sein können, nur teilweise bestimmte Antworten gibt. So verhält es sich bei der Bestimmung des Wertes des Integrals $\iint y\,dr\,d\varphi$. Soll es allgemein von $r=k$ bis $r=l$ und von $\varphi=K$ bis $\varphi=L$ genommen werden und wird der Wert von $\frac{u}{t}$
\[
\begin{array}{rcl}
\hbox{für }r=k,\quad\varphi=K&\hbox{durch}&\theta\\
r=k,\quad\varphi=L&\hbox{durch}&\theta'\\
r=l,\quad\varphi=K&\hbox{durch}&\theta''\\
r=l,\quad\varphi=L&\hbox{durch}&\theta'''
\end{array}
\]
bezeichnet, so ergeben die analytischen Operationen für jenes Integral den Ausdruck
\[
\mathrm{Arc\,tg}\,\theta-\mathrm{Arc\,tg}\,\theta'
-\mathrm{Arc\,tg}\,\theta''+\mathrm{Arc\,tg}\,\theta'''
\]

Das Integral hat tatsächlich nur dann einen wahren Wert, wenn $y$ innerhalb der angegebenen Grenzen stets endlich bleibt. Dieser wahre Wert ist dann zwar im obigen Ausdruck enthalten, aber durch ihn allein noch nicht vollständig bestimmt, da die Funktion Arc. tg. mehrwertig ist. Erst durch weitere, im Übrigen nicht schwierige Betrachtungen muss entschieden werden, \emph{welche}
\SourceLeaf{121}
Werte dieser Funktion im jeweiligen Fall zu nehmen sind. Wird dagegen $y$ irgendwo innerhalb der angegebenen Grenzen unendlich, so ist die Frage nach dem Wert von $\iint y\,dr\,d\varphi$ eigentlich unsinnig. Wird dennoch der Versuch unternommen, sie unter Missachtung dieses Unsinns zu beantworten, so darf es nicht befremden, auf einem Weg diese und auf anderen Wegen jene Antwort zu erhalten. Diese verschiedenen Antworten sind jedoch stets im obigen teilweise bestimmten Ausdruck enthalten.

\par\medskip\noindent\rule{0.25\linewidth}{0.4pt}\par\smallskip
\begingroup\small
\emph{Handschriftliche Bemerkungen.}

Die \emph{Bedingung} für $R$ lässt sich auch so formulieren, dass es größer ist als $1$ und als die absolute Summe der Größen $A,B,C$ usw., multipliziert mit $\surd2$. Diese absolute Summe sei $=\alpha A+\beta B+\gamma C+\delta D+\mathrm{usw.}$, sodass jeder der Koeffizienten $\alpha,\beta,\gamma$ usw. entweder $+1$ oder $-1$ ist. Dann ist der positive Wert von $T$ usw. aus folgender Form ersichtlich:
\[
\begin{aligned}&T=R^{m-1}\cos45^0\bigl(R-\surd2.(\alpha A+\beta B+\gamma C+\mathrm{usw.})\bigr)\\
&\quad+R^{m-1}\bigl(\alpha A+A\cos(45^0+\varphi)\bigr)\\
&\quad+R^{m-2}\bigl(\beta BR+B\cos(45^0+2\varphi)\bigr)+\mathrm{usw.}
\end{aligned}
\]

\emph{Satz.} Sind $a,b,c\,.\,.\,.\,m,n$ die Wurzeln der Gleichung $fx=0$ und $a',b',c'\,.\,.\,.\,m'$ die Wurzeln der Gleichung $f'x=0$, wobei $f'x=\frac{dfx}{dx}$ gilt, und bezeichnen dieselben Buchstaben die entsprechenden Punkte in der Ebene, so herrscht, wenn in $a,b,c\,.\,.\,.\,m,n$ gleiche abstoßende oder anziehende Massen gedacht werden, die im umgekehrten Verhältnis der Entfernung wirken, in $a',b',c'\,.\,.\,.\,m'$ Gleichgewicht.
\endgroup

\par\medskip\begin{center}\rule{0.25\linewidth}{0.4pt}\end{center}

